% Continuation of the single canonical lesson, not a second theory pass.
\subsection{Framing Efficiency: কার্যকর ডেটার অনুপাত}
\label{ch2:c02-07}
\begin{keypointbox}{কতটুকু কার্যকর ডেটা গেল?}
পাঠানো মোট বিটের মধ্যে কার্যকর ডেটা বিটের শতকরা অনুপাতকে ফ্রেমিং দক্ষতা (Framing Efficiency) বলে।
\[
\eta=\frac{\text{কার্যকর Data Bit}}{\text{মোট প্রেরিত Bit}}\times100\%
\]
অ্যাসিনক্রোনাস পদ্ধতিতে $N$টি ক্যারেক্টার পাঠানো হলে এবং প্রতিটিতে $d$ ডেটা বিট থাকলে, $s,p,t$ যথাক্রমে স্টার্ট, প্যারিটি ও স্টপ বিট বোঝায়। মোট প্রেরিত বিট হয় $N(d+s+p+t)$। সিনক্রোনাস ব্লকে $Nd$ ডেটা বিট ও মোট $H$ বিট বাড়তি তথ্য থাকলে মোট বিট $Nd+H$। এই সরল হিসাবে বিরতি ও প্রোটোকলের অতিরিক্ত তথ্য ধরা হয়নি। বাইট থেকে বিটে রূপান্তরের নিয়ম অংশ \ref{ch2:c01-09}-এ আছে।
\end{keypointbox}
\chTwoPair{\textbf{উদাহরণ ১: ক্যারেক্টারের ফ্রেমের বিন্যাস বদলালে}\par
8 Data + 1 Start + 1 Stop হলে $\eta=8/10\times100=80\%$। একই ফ্রেমের বিন্যাসে বিরতিহীনভাবে 1000 bps পাঠালে কার্যকর ডেটার হার $1000\times0.8=800$ bps।\par
স্টপ বিট 2টি হলে $8/11\times100\approx72.73\%$। এর সঙ্গে 1 প্যারিটি বিট থাকলে $8/12\times100\approx66.67\%$। তাই হিসাবের আগে ফ্রেমের বিন্যাস বুঝে নিতে হবে।}{\textbf{উদাহরণ ২: সেন্সরের সংযোগ বাছাই}\par
সেন্সর প্রতি সেকেন্ডে 120টি ক্যারেক্টার পাঠায়। প্রতিটিতে 8 ডেটা, 1 স্টার্ট ও 2 স্টপ বিট থাকে। দরকার $120\times11=1320$ bps; শুধু কার্যকর ডেটা $120\times8=960$ bps।\par
960 bps সংযোগ শুধু মূল ডেটা বা পে-লোড (Payload) বহনের চাহিদা মেটায়; পুরো ফ্রেম পাঠাতে যথেষ্ট নয়। অন্তত 1320 bps এবং বাস্তব ব্যবস্থার অতিরিক্ত চাহিদা বিবেচনা করতে হবে।}
\begin{examplebox}{উদাহরণ ৩: একই ডেটা, ক্যারেক্টার বনাম ব্লক}
\textbf{নির্দিষ্ট ফ্রেমের বিন্যাস:} 1000 Byte ডেটা; প্রতি Byte-এ 8 ডেটা বিট। এই হিসাবে হেডার ও ট্রেলারের নির্দিষ্ট দৈর্ঘ্য ধরা হয়েছে। সব প্রোটোকলের ফ্রেমের বিন্যাস একই নয়।
\chTwoPair{\textbf{Asynchronous}\par
ক্যারেক্টারপ্রতি 1 স্টার্ট ও 2 স্টপ বিট; তাই মোট $1+8+2=11$ বিট।\par
কার্যকর $1000\times8=8000$ বিট; মোট $1000\times11=11000$ বিট।\par
ফ্রেমিং দক্ষতা $=8000/11000\times100\approx72.73\%$।}{\textbf{Synchronous}\par
হেডার 2 Byte ও ট্রেলার 2 Byte; বাড়তি তথ্য $4\times8=32$ বিট।\par
মোট $8000+32=8032$ বিট।\par
ফ্রেমিং দক্ষতা $=8000/8032\times100\approx99.60\%$।}
এই উদাহরণের হেডারে 1 Byte ফ্ল্যাগ ও 1 Byte নিয়ন্ত্রণক্ষেত্র থাকতে পারে। ট্রেলার ব্লকের শেষ ও যাচাইয়ের তথ্য দেয়। একই সংযোগের হার এবং বিরতিহীন প্রেরণ ধরে দ্বিতীয় পদ্ধতিতে কম বিট পাঠাতে হয়। 100টি 8-বিট ক্যারেক্টারের ব্লকে বাড়তি তথ্য 40 বিট হলে $800/840\times100\approx95.24\%$; ব্লকের আকার বা বাড়তি তথ্য বদলালে ফলও বদলায়। বড় ব্লক ব্যবহারে বাফারের প্রয়োজন ও ত্রুটির পরে পুনঃপ্রেরণের খরচ বিবেচনা করতে হয়।
\end{examplebox}

\subsection{দিকভিত্তিক Mode: Simplex}
\label{ch2:c02-08}
\begin{definitionbox}{Simplex Mode}
ডেটা শুধু একদিকে পাঠানো যায়—এমন যোগাযোগকে সিমপ্লেক্স মোড (Simplex Mode) বলে। এক পক্ষ প্রেরক এবং অন্য পক্ষ প্রাপক হিসেবে কাজ করে।
\end{definitionbox}
টেলিভিশন ও রেডিও সম্প্রচারে অনুষ্ঠান প্রেরক থেকে দর্শক বা শ্রোতার কাছে যায়। কিবোর্ডের ইনপুট কম্পিউটারে পাঠানো এবং কম্পিউটার থেকে মনিটরে প্রদর্শনের ডেটা পাঠানো একমুখী প্রবাহের উদাহরণ। আধুনিক যন্ত্রের ভৌত সংযোগে নিয়ন্ত্রণ বা অবস্থার তথ্য ফিরতে পারে। তাই উদাহরণে কোন ডেটার প্রবাহ বোঝানো হয়েছে, তা বুঝতে হবে।

\subsection{Half-duplex: দুইদিকে, পালাক্রমে}
\label{ch2:c02-09}
\begin{definitionbox}{Half-duplex Mode}
দুই পক্ষই ডেটা পাঠাতে ও গ্রহণ করতে পারে, তবে একই সময়ে পারে না—এমন যোগাযোগকে হাফ-ডুপ্লেক্স মোড (Half-duplex Mode) বলে। এক পক্ষ পাঠালে অন্য পক্ষ গ্রহণ করে; পরের পালায় তাদের ভূমিকা বদলায়।
\end{definitionbox}
ওয়াকিটকিতে বোতাম চেপে একজন কথা বললে অন্যজন শোনেন। দুইজন একসঙ্গে কথা পাঠাতে পারেন না। একই পথ পালাক্রমে ব্যবহার করা যায়, তবে উত্তর দিতে নিজের পালার জন্য অপেক্ষা করতে হয়।
\begin{boardbox}{SMS/FAX-এর প্রচলিত উদাহরণ কীভাবে পড়বে?}
SMS-এর পূর্ণরূপ Short Message Service; FAX শব্দটি Facsimile থেকে এসেছে। পরীক্ষার কিছু প্রচলিত উদাহরণে SMS/FAX-এর বার্তা ও ফিরতি বার্তা পাঠানোকে পালাক্রমিক যোগাযোগ ধরে হাফ-ডুপ্লেক্স দেখানো হয়। প্রশ্নে এই প্রসঙ্গ থাকলে তা বিবেচনা করতে হবে। তবে মানুষ পালাক্রমে SMS বা চ্যাটের উত্তর দেন বলেই মোবাইল নেটওয়ার্কের ভৌত সংযোগ হাফ-ডুপ্লেক্স হয় না। সংযোগটি একই সময়ে দুইদিকে ডেটা পাঠাতে পারে কি না, সেটিই নির্ধারক।
\end{boardbox}

\subsection{Full-duplex: একই সময়ে দুইদিকে}
\label{ch2:c02-10}
\begin{definitionbox}{Full-duplex Mode}
দুই পক্ষ একই সময়ে ডেটা পাঠাতে ও গ্রহণ করতে পারে—এমন যোগাযোগকে ফুল-ডুপ্লেক্স মোড (Full-duplex Mode) বলে।
\end{definitionbox}
ফোন বা ভিডিও কলে একই সময়ে কথা বলা ও শোনা যায়। পারস্পরিক যোগাযোগে পালা বদলের অপেক্ষা করতে হয় না। তবে বেশি ডেটা হার বা সব কাজে বেশি সুবিধা ফুল-ডুপ্লেক্সের সংজ্ঞার অংশ নয়। একমুখী ঘোষণায় সিমপ্লেক্সই যথেষ্ট হতে পারে। সংযোগের দিক জানতে শুধু যন্ত্রের নাম বা সময় ভাগ করে ব্যবহারের পদ্ধতির নাম যথেষ্ট নয়।
\chTwoPair{\centering\begin{tikzpicture}[x=1cm,y=1cm]
\foreach \y/\lab in {0/Simplex,-1.6/{Half duplex\\পালাক্রমে},-3.2/{Full duplex\\একই সময়ে}}{
\node[chTwoNode] at (0,\y){A};
\node[chTwoNode] at (5.5,\y){B};
\node[font=\normalsize,align=center] at (2.75,\y+.85){\lab};}
\draw[chTwoArrow] (.4,0)--(5.1,0);
\draw[chTwoArrow] (.4,-1.45)--(5.1,-1.45);
\draw[chTwoArrow,dashed] (5.1,-1.75)--(.4,-1.75);
\draw[chTwoArrow] (.4,-3.05)--(5.1,-3.05);
\draw[chTwoArrow] (5.1,-3.35)--(.4,-3.35);
\end{tikzpicture}}{\textbf{চিত্র পড়ার নিয়ম}\par
Simplex: শুধু $A\rightarrow B$।\par
Half-duplex: প্রথম পালায় $A\rightarrow B$, পরে $B\rightarrow A$।\par
Full-duplex: একই সময়ে $A\rightarrow B$ এবং $B\rightarrow A$।\par
\textbf{বাস্তব সিদ্ধান্ত:} পাহারাদারের সংক্ষিপ্ত পালাক্রমিক কথায় Walkie-talkie চলে। দীর্ঘ পারস্পরিক ভিডিও আলোচনায় একই সময়ে কথা বলা ও শোনার সুবিধা দরকার। প্রয়োজনের সঙ্গে খরচ ও ব্যবস্থাও মিলিয়ে নেবে।}

\subsection{প্রাপকভিত্তিক Delivery: Unicast}
\label{ch2:c02-11}
\begin{definitionbox}{Unicast}
এক প্রেরক থেকে একজন নির্দিষ্ট প্রাপকের কাছে ডেটা পাঠানোর পদ্ধতিকে ইউনিকাস্ট (Unicast) বলে।
\end{definitionbox}
এক বন্ধুকে ই-মেইল বা SMS পাঠানো এবং একজন ব্যবহারকারীর নির্দিষ্ট সার্ভারে অনুরোধ পাঠানো ইউনিকাস্টের উদাহরণ। এর অর্থ একজন থেকে একজনের কাছে পাঠানো বা One-to-one। তবে ওই সংযোগ একমুখী না একই সময়ে দুইমুখী, তা ইউনিকাস্ট নাম থেকে জানা যায় না।

\subsection{Broadcast: নির্ধারিত পরিসরের সবাই}
\label{ch2:c02-12}
\begin{definitionbox}{Broadcast}
এক প্রেরক থেকে নির্দিষ্ট সম্প্রচার পরিসর বা নেটওয়ার্কের ব্রডকাস্ট ডোমেইনের সব প্রাপকের উদ্দেশে ডেটা পাঠানোর পদ্ধতিকে ব্রডকাস্ট (Broadcast) বলে।
\end{definitionbox}
টেলিভিশন বা রেডিও সম্প্রচার এবং স্থানীয় নেটওয়ার্কের ব্রডকাস্ট বার্তা এর উদাহরণ। এখানে সবাই বলতে নির্দিষ্ট পরিসরের সবাইকে বোঝায়; ইন্টারনেটের প্রত্যেক যন্ত্রকে নয়। ব্রডকাস্ট প্রাপকের ভিত্তিতে এবং সিমপ্লেক্স দিকের ভিত্তিতে নির্ধারিত হয়। একটি টেলিভিশন সম্প্রচার একই সঙ্গে দুটির উদাহরণ হতে পারে। রেডিও তরঙ্গ মাধ্যম বোঝায়, আর ব্রডকাস্ট বোঝায় কারা বার্তাটি পাবেন।

\subsection{Multicast: নির্বাচিত প্রাপকের দল}
\label{ch2:c02-13}
\begin{definitionbox}{Multicast}
এক প্রেরক থেকে নির্বাচিত প্রাপক দলের উদ্দেশে ডেটা পাঠানোর পদ্ধতিকে মাল্টিকাস্ট (Multicast) বলে। IP Multicast-এ দলের সদস্যদের কাছে ডেটা পাঠাতে একটি দলীয় ঠিকানা বা Group Address ব্যবহার করা হয়।
\end{definitionbox}
মাল্টিকাস্ট সমর্থন করে এমন IPTV ব্যবস্থায় নিবন্ধিত দর্শকদের একই ডেটা প্রবাহ দেওয়া যায়। সবার কাছে অপ্রয়োজনীয় ডেটা পাঠানোর বদলে নির্বাচিত দলকে পাঠানো হয়। ব্রডকাস্টে নির্দিষ্ট পরিসরের সবাই প্রাপক; মাল্টিকাস্টে প্রাপক একটি নির্বাচিত দল।
\begin{mistakebox}{দলগত বার্তা মানেই IP Multicast নয়}
পরীক্ষায় একজন থেকে নির্বাচিত দলের কাছে পাঠানোর প্রসঙ্গে Group SMS বা Chatting-কে মাল্টিকাস্ট বলা হতে পারে। বাস্তবে একই বার্তার পৃথক ইউনিকাস্ট কপিও দলের সদস্যদের পাঠানো যায়; একে Multiple Unicast বলে। তাই প্রশ্নে শুধু দলে পাঠানো, নাকি দলীয় ঠিকানায় IP Multicast বোঝানো হয়েছে, তা বুঝতে হবে। ভিডিও কনফারেন্সের প্রাপক সংখ্যা ও ব্যবস্থার তথ্য না জেনে তাকে নিশ্চিত ইউনিকাস্ট বা মাল্টিকাস্ট বলা যায় না।
\end{mistakebox}
\par\noindent\begin{minipage}{\linewidth}\centering
\begin{tikzpicture}[x=1cm,y=1cm]
\foreach \offset/\lab/\id in {0/Unicast/u,5.4/Broadcast/b,10.8/Multicast/m}{
\node[font=\normalsize] at (\offset+2,1.6){\lab};
\node[chTwoNode,font=\normalsize] (s\id) at (\offset,0){S};
\foreach \num/\y in {1/1,2/0,3/-1}{\node[chTwoNode,font=\normalsize] (r\id\num) at (\offset+3.8,\y){R\num};}}
\draw[chTwoArrow] (su)--(ru2);
\foreach \x in {1,2,3}{\draw[chTwoArrow] (sb)--(rb\x);}
\foreach \x in {1,3}{\draw[chTwoArrow] (sm)--(rm\x);}
\end{tikzpicture}\end{minipage}\par
\begin{tabularx}{\linewidth}{|p{.18\linewidth}|X|X|}\hline
\textbf{প্রাপকভিত্তিক নাম} & \textbf{প্রাপকসংখ্যার বর্ণনা} & \textbf{চিত্রের প্রাপক}\\\hline
Unicast & Point-to-point: একজন থেকে একজন & শুধু R2\\\hline
Broadcast & Point-to-multipoint: একজন থেকে ওই পরিসরের সবাই & R1, R2 ও R3\\\hline
Multicast & নির্বাচিত Multipoint: নির্দিষ্ট দল & দলের R1 ও R3; R2 নয়\\\hline
\end{tabularx}
\textbf{নামের প্রসঙ্গ:} এই ছকে Point শব্দযুক্ত নামগুলো প্রাপকের সংখ্যা বোঝায়। পরের টপিকে ভৌত Point-to-point ও Multipoint দিয়ে সংযোগের বিন্যাস বোঝানো হয়েছে। মাল্টিপয়েন্ট ব্যবস্থাতেও একজন প্রাপককে ইউনিকাস্ট ডেটা দেওয়া যায়। তাই সংযোগের নাম থেকে ব্রডকাস্ট বা মাল্টিকাস্ট নিশ্চিত করা যায় না।
ব্যক্তিগত বার্তায় ইউনিকাস্ট, সবার উদ্দেশে ঘোষণায় ব্রডকাস্ট এবং সমর্থিত নেটওয়ার্কে দলের ডেটা প্রবাহে মাল্টিকাস্ট উপযোগী। দল কত বড়, ব্যবস্থা কোন পদ্ধতি সমর্থন করে এবং একই ডেটা কতবার পাঠাতে হয়—এসব বিষয় বিবেচনা করতে হবে।
\chTwoPair{\textbf{আরও পরিচয়—Anycast:} একই সেবার ঠিকানা থাকা একাধিক সম্ভাব্য প্রাপকের মধ্যে রাউটিংয়ের হিসাবে কাছের একটির কাছে ডেটা পাঠানো হয়। যেমন, একাধিক স্থানে থাকা DNS সার্ভারের একটি থেকে সেবা নেওয়া। এখানে কাছের অর্থ রাউটিং দূরত্বে কাছে; শুধু মানচিত্রের দূরত্ব নয়।}{\textbf{Geocast:} নির্দিষ্ট ভৌগোলিক এলাকার প্রাপকদের উদ্দেশে বার্তা পাঠানোর ধারণা। উপকূলের নির্দিষ্ট অঞ্চলের বাসিন্দাদের ঝড়ের সতর্কতা দেওয়ার প্রয়োজন এর উদাহরণ। প্রাপকদের অবস্থান দিয়ে দল নির্বাচন করা হয়। তবে দলের সব সদস্যের কাছে একই IP Multicast পদ্ধতিতে বার্তা পাঠানো হবে, এমন নিশ্চয়তা নেই।}
